\chapter{$\En$ Koszul duality and higher-dimensional bar complexes}
\label{ch:en-koszul-duality}
\index{En algebra@$\En$ algebra!Koszul duality|textbf}
\index{bar complex!higher-dimensional|textbf}

Throughout this monograph, the bar differential has been built from
configuration spaces of points on a \emph{curve}: the propagator
$\omega_{ij} = d\log(z_i - z_j)$ is a meromorphic $1$-form, its
singularities are controlled by Arnold relations, and the bar-cobar
adjunction lives over the chiral operad.  What changes when the
curve is replaced by an $n$-manifold?  The propagator becomes
a closed $(n-1)$-form, singularities along collision strata are
controlled by \emph{Totaro relations}, and the resulting bar-cobar
adjunction governs $\En$-algebras.  This chapter develops the
general theory, proves the $n = 1$ and $n = 2$ cases in detail,
and formulates the programme for arbitrary~$n$.

The passage from $n = 1$ to general~$n$ is not an idle
generalization.  The $n = 2$ bar complex governs
$\Etwo$-algebras---the natural algebraic structure on braided
monoidal categories---and connects to Kontsevich's formality theorem
for the little $2$-disks operad.  The $n = 3$ case interfaces with
Chern--Simons theory and Vassiliev invariants
(Chapter~\ref{ch:kontsevich-integral}).  Each step up the
dimensional ladder trades holomorphic structure for topological
generality, with Arnold relations ($n = 1$) specializing Totaro
relations ($n \geq 2$).


%================================================================
% SECTION 1: HIGHER-DIMENSIONAL CONFIGURATION SPACES
%================================================================

\section{Higher-dimensional configuration spaces}
\label{sec:higher-dim-config}
\index{configuration space!higher-dimensional}

\begin{definition}[Configuration space of $\bR^n$]
\label{def:config-rn}
For a smooth oriented $n$-manifold~$M$, the \emph{ordered
configuration space of $k$~points} is
\[
\Conf_k(M) \;=\;
\bigl\{(x_1, \ldots, x_k) \in M^k
  \mid x_i \neq x_j \text{ for } i \neq j\bigr\}.
\]
For $M = \bR^n$, the space $\Conf_k(\bR^n)$ is a smooth manifold
of real dimension~$nk$, and its cohomology admits an explicit
generators-and-relations presentation due to Arnold ($n = 2$, i.e., $\bR^n = \bC$) and
Totaro (general $n$).
\end{definition}

\begin{theorem}[Arnold presentation {\cite{Arnold69}}; $\bC \cong \bR^2$;
\ClaimStatusProvedElsewhere]
\label{thm:arnold-presentation}
\index{Arnold relations!presentation}
The cohomology ring $H^*(\Conf_k(\bC); \bQ)$ is the graded-commutative
algebra generated by classes $\omega_{ij} \in H^1$,
$1 \leq i < j \leq k$, subject to the \emph{Arnold relations}:
\begin{equation}\label{eq:arnold-rel-en}
\omega_{ij} \wedge \omega_{jk}
+ \omega_{jk} \wedge \omega_{ki}
+ \omega_{ki} \wedge \omega_{ij} \;=\; 0
\quad (i < j < k).
\end{equation}
This is the Orlik--Solomon algebra $\mathrm{OS}_k$
\textup{(}Theorem~\textup{\ref{thm:os-cohomology-config})}.
\end{theorem}

\begin{theorem}[Totaro presentation, general $n$ {\cite{Coh76}};
\ClaimStatusProvedElsewhere]
\label{thm:totaro-presentation}
\index{Totaro relations|textbf}
\index{configuration space!cohomology ring}
For $n \geq 2$, the cohomology ring
$H^*(\Conf_k(\bR^n); \bQ)$ is the graded-commutative algebra
generated by classes $\alpha_{ij} \in H^{n-1}$,
$1 \leq i \neq j \leq k$, subject to the relations:
\begin{enumerate}[label=\textup{(\roman*)}]
\item \textup{(Symmetry)}\quad
  $\alpha_{ij} = (-1)^n \alpha_{ji}$;
\item \textup{(Totaro relations)}\quad
  $\alpha_{ij} \cdot \alpha_{jk}
  + \alpha_{jk} \cdot \alpha_{ki}
  + \alpha_{ki} \cdot \alpha_{ij} = 0$
  for distinct $i, j, k$;
\item \textup{(Degree)}\quad
  $\alpha_{ij}^2 = 0$ when $n$ is odd;
  when $n$ is even, $\alpha_{ij}^2$ is determined by the Euler class.
\end{enumerate}
For $n = 1$ and ordered $i < j$, setting $\alpha_{ij} = \omega_{ij}$
recovers the Arnold presentation.  For $n = 2$, the generators
$\alpha_{ij} \in H^1$ satisfy the same formal relations as
Arnold's generators, but with the additional symmetry
$\alpha_{ij} = \alpha_{ji}$ (since $(-1)^2 = 1$).
\end{theorem}

\begin{proof}
See Totaro~\cite{Totaro96}, Theorem~1; also
Cohen~\cite{Cohen76}, Theorem~A for the formality statement.
The key geometric input is that projection maps
$\Conf_k(\bR^n) \to \Conf_2(\bR^n)$ induce pullbacks of the
fundamental class of $S^{n-1}$ (the linking sphere), and all
relations among these pullbacks are generated by the three-term
relation (ii).
\end{proof}

\begin{remark}[Arnold as Totaro at $n = 2$]
\label{rem:arnold-totaro}
\index{Arnold relations!as Totaro specialization}
In the $\En$ hierarchy, $n$ always refers to the \emph{real} dimension.
At $n = 1$, the generators $\alpha_{ij} \in H^0(\Conf_2(\bR))$ would
live in degree~$0$, which is the associative (non-chiral) setting.
The classical Arnold generators on $\Conf_k(\bC)$ are Totaro
generators at $n = 2$ (real dimension of $\bC \cong \bR^2$).
Our chiral bar complex on algebraic curves uses $n = 2$
(complex dimension~$1$, real dimension~$2$).  Thus:
\begin{center}
\begin{tabular}{ccc}
$\En$-level & Manifold & Bar complex \\
\hline
$n = 1$ & interval/circle & associative bar (classical) \\
$n = 2$ & surface/curve & chiral bar (this monograph) \\
$n = 3$ & $3$-manifold & Chern--Simons bar \\
general $n$ & $n$-manifold & $\En$ bar
\end{tabular}
\end{center}
Our Main Theorems~A, B, C are the $n = 2$ case of $\En$ Koszul
duality.
\end{remark}


\subsection{Fulton--MacPherson compactification in dimension $n$}
\label{subsec:fm-general-n}
\index{Fulton--MacPherson compactification!higher-dimensional}

\begin{definition}[FM compactification of $\Conf_k(\bR^n)$]
\label{def:fm-rn}
The \emph{Fulton--MacPherson compactification}
$\overline{\Conf}_k(\bR^n)$ is the real oriented blowup of $(\bR^n)^k$
along all partial diagonals, ordered by reverse inclusion.
Equivalently (Axelrod--Singer, Kontsevich), it is the closure of
the image of the map
\[
\Conf_k(\bR^n)
\;\hookrightarrow\;
(\bR^n)^k
\times \prod_{i<j} S^{n-1}
\times \prod_{|S| \geq 2} [0, \infty]
\]
recording positions, relative directions
$\theta_{ij} = (x_i - x_j)/|x_i - x_j| \in S^{n-1}$,
and relative distances.  The result is a smooth manifold with
corners whose boundary strata are indexed by nested partitions
of $\{1, \ldots, k\}$---the same combinatorial data as the
Stasheff associahedron ($n = 1$) or the FM operad ($n \geq 2$).
\end{definition}

\begin{proposition}[Boundary strata and operadic structure;
\ClaimStatusProvedElsewhere{} \cite{FM94}]
\label{prop:fm-boundary-strata}
\index{Fulton--MacPherson!boundary strata}
The codimension-$1$ boundary strata of
$\overline{\Conf}_k(\bR^n)$ are indexed by subsets
$S \subset \{1, \ldots, k\}$ with $|S| \geq 2$, and are
diffeomorphic to
\[
\partial_S \overline{\Conf}_k(\bR^n)
\;\cong\;
\overline{\Conf}_{k-|S|+1}(\bR^n)
\times
\overline{\Conf}_{|S|}(\bR^n).
\]
This factorization equips $\{\overline{\Conf}_k(\bR^n)\}_{k \geq 0}$
with the structure of a topological operad---the
\emph{Fulton--MacPherson operad} $\mathsf{FM}_n$---which is a model
for the $\En$ operad \textup{(}i.e., weakly equivalent to the
little $n$-disks operad\textup{)}.
\end{proposition}

\begin{proof}
The operadic structure is due to Getzler--Jones~\cite{GJ}
and Salvatore; the equivalence $\mathsf{FM}_n \simeq \En$ is
proved by Salvatore--Longoni for $n \geq 2$ and follows from
Stasheff's theorem for $n = 1$.
\end{proof}


%================================================================
% SECTION 2: THE E_n PROPAGATOR
%================================================================

\section{The $\En$ propagator}
\label{sec:en-propagator}
\index{propagator!$\En$}

\begin{definition}[$\En$ propagator]
\label{def:en-propagator}
\index{propagator!Green's form}
An \emph{$\En$ propagator} on an oriented $n$-manifold~$M$ is a
smooth closed $(n-1)$-form
\[
G \;\in\; \Omega^{n-1}(M \times M \setminus \Delta)
\]
satisfying the following properties:
\begin{enumerate}[label=\textup{(\roman*)}]
\item \textup{(Poincar\'e dual of diagonal)}\quad
  $G$ represents the Poincar\'e dual of the diagonal
  $\Delta \subset M \times M$ in the sense that for any
  closed form $\eta \in \Omega^*(M)$,
  \[
  \int_{M \times M} G \wedge \pi_2^* \eta
  \;=\; \int_M \eta,
  \]
  where $\pi_2\colon M \times M \to M$ is the second projection.
\item \textup{(Singularity)}\quad
  Near the diagonal $\Delta$, in geodesic coordinates,
  \[
  G \;\sim\;
  \frac{1}{\mathrm{vol}(S^{n-1})}\;
  \frac{r_1\, dr_2 \wedge \cdots \wedge dr_n
  - r_2\, dr_1 \wedge dr_3 \wedge \cdots \wedge dr_n
  + \cdots}{|r|^n}
  \;+\; \text{smooth},
  \]
  where $r = x_1 - x_2$ and $\mathrm{vol}(S^{n-1}) = 2\pi^{n/2}/\Gamma(n/2)$.
\item \textup{(Symmetry)}\quad
  Under the involution $\sigma\colon (x_1, x_2) \mapsto (x_2, x_1)$,
  $\sigma^* G = (-1)^n G$.
\end{enumerate}
\end{definition}

\begin{example}[$n = 1$: the chiral propagator]
\label{ex:n1-propagator}
For $M = \bC$ (real dimension $2$, but complex dimension $1$---the
chiral case), the propagator is
$G = \omega_{12} = d\log(z_1 - z_2)$, a meromorphic $1$-form.
It is closed ($dG = 0$ away from $z_1 = z_2$), represents the
Poincar\'e dual of the diagonal, and satisfies
$\sigma^* G = d\log(z_2 - z_1) = d\log(-(z_1-z_2)) = d\log(-1) + d\log(z_1-z_2) = G$,
since $d\log(-1) = 0$ (differential of a constant).  Residues along $z_1 = z_2$ extract OPE data.
\end{example}

\begin{example}[$n = 2$: the $\Etwo$ propagator on $\bR^2$]
\label{ex:n2-propagator}
For $M = \bR^2$, the propagator is
\[
G \;=\;
\frac{1}{2\pi}\;
\frac{(x_1 - x_2)\, dy_1 - (y_1 - y_2)\, dx_1}
     {(x_1 - x_2)^2 + (y_1 - y_2)^2}
\;=\;
\frac{1}{2\pi}\, d\arg(z_1 - z_2),
\]
a closed $1$-form on $\Conf_2(\bR^2)$.  This is the angular form
(winding number kernel): integration over a small circle
$S^1 \subset \bR^2$ around $x_2$ gives $1$.  Note that
$G = (1/2\pi)\, d\arg(z_1 - z_2)$ is exactly the
\emph{Kontsevich propagator} on $S^1 \subset \bC$
(Proposition~\ref{prop:propagator-restriction}).  The $\Etwo$
propagator on a surface and the chiral propagator on a complex curve
are thus two aspects of the same object: the former is real-valued
and topological, the latter is holomorphic and algebraic.
\end{example}

\begin{example}[$n = 3$: the Chern--Simons propagator]
\label{ex:n3-propagator}
\index{Chern--Simons!propagator}
For $M = \bR^3$, the propagator is
\[
G \;=\;
\frac{1}{4\pi}\;
\frac{\epsilon_{abc}\, r^a\, dr^b \wedge dr^c}{|r|^3},
\]
where $r = x_1 - x_2 \in \bR^3$.  This is the closed $2$-form on
$\Conf_2(\bR^3)$ representing the linking number:
$\int_{S^2} G = 1$.  In the Axelrod--Singer perturbative expansion of
Chern--Simons theory~\cite{Kon94}, this propagator appears as the
kernel of the gauge-fixed Green's operator $d^{-1}$.
\end{example}

\begin{proposition}[Residue as linking sphere integral;
\ClaimStatusProvedHere]
\label{prop:linking-sphere-residue}
\index{residue!linking sphere}
For an $\En$ propagator $G$ on $\bR^n$ and a smooth function
$f$ on a neighborhood of the diagonal, the \emph{$\En$ residue} is:
\begin{equation}\label{eq:en-residue}
\Res_\Delta(G \wedge f)
\;:=\;
\lim_{\varepsilon \to 0}
\int_{|x_1 - x_2| = \varepsilon} G \cdot f(x_1, x_2)
\;=\; f(x, x),
\end{equation}
where the integral is over the linking sphere
$S^{n-1}_\varepsilon \subset \Conf_2(\bR^n)$ of radius~$\varepsilon$
around the diagonal.  This generalizes the residue operation
$\Res_{z_1 = z_2}$ in the chiral ($n = 2$) case.
\end{proposition}

\begin{proof}
By the singularity condition (ii) of
Definition~\ref{def:en-propagator}, $G$ restricts on the
sphere $|r| = \varepsilon$ to the standard volume form on $S^{n-1}$
(normalized to total volume~$1$) plus terms vanishing as
$\varepsilon \to 0$.  Hence the integral extracts the value
$f(x, x)$, exactly as a contour integral residue extracts the
leading Laurent coefficient.
\end{proof}


%================================================================
% SECTION 3: THE E_2 BAR COMPLEX IN DETAIL
%================================================================

\section{The $\Etwo$ bar complex}
\label{sec:e2-bar}
\index{bar complex!$\Etwo$}

We now construct the $\En$ bar complex in the case $n = 2$,
connecting it to the chiral bar-cobar duality of
Chapter~\ref{chap:bar-cobar}.

\begin{definition}[$\Etwo$ bar complex]
\label{def:e2-bar-complex}
\index{bar complex!$\Etwo$|textbf}
Let $A$ be an $\Etwo$-algebra (equivalently, an algebra over
the little $2$-disks operad).  The \emph{$\Etwo$ bar complex}
is the chain complex
\[
\barB_{\Etwo}(A)
\;=\;
\bigoplus_{k \geq 1}
\Omega^*\bigl(\overline{\Conf}_k(\bR^2)\bigr)
\otimes_{\Sigma_k} (s^{-1} \bar{A})^{\otimes k},
\]
with differential
\begin{equation}\label{eq:e2-bar-diff}
d_{\Etwo}
\;=\;
d_{\mathrm{dR}}
\;+\;
\sum_{S \subset \{1, \ldots, k\},\, |S| \geq 2}
\Res_{\partial_S}\bigl(G_S \wedge \mu_S\bigr),
\end{equation}
where:
\begin{itemize}
\item $G_S = \bigwedge_{i < j,\; i, j \in S} \alpha_{ij}$ is the
  product of propagators for pairs within the
  colliding subset~$S$,
\item $\Res_{\partial_S}$ is the linking-sphere residue along the
  boundary stratum $\partial_S$ of
  $\overline{\Conf}_k(\bR^2)$
  \textup{(}Proposition~\textup{\ref{prop:linking-sphere-residue})},
\item $\mu_S\colon \bar{A}^{\otimes |S|} \to \bar{A}$ is the
  $\Etwo$-multiplication indexed by the FM operad element
  corresponding to the stratum~$\partial_S$.
\end{itemize}
\end{definition}

\begin{theorem}[$d^2 = 0$ from Totaro relations;
\ClaimStatusProvedHere]
\label{thm:e2-d-squared}
\index{Totaro relations!and $d^2 = 0$}
The differential $d_{\Etwo}$ of
Definition~\textup{\ref{def:e2-bar-complex}} satisfies
$d_{\Etwo}^2 = 0$.
\end{theorem}

\begin{proof}
The proof follows the same structure as the chiral case
(Corollary~\ref{cor:bar-d-squared-zero-arnold}), with the Arnold
relations replaced by Totaro relations.

We compute $d_{\Etwo}^2$ on a bar element
$\alpha \in \Omega^*(\overline{\Conf}_k(\bR^2))
\otimes (s^{-1} \bar{A})^{\otimes k}$.
The square decomposes into three terms:
\[
d_{\Etwo}^2
\;=\;
\underbrace{d_{\mathrm{dR}}^2}_{= 0}
\;+\;
\underbrace{[d_{\mathrm{dR}},\, d_{\mathrm{res}}]}_{(\star)}
\;+\;
\underbrace{d_{\mathrm{res}}^2}_{(\star\star)},
\]
where $d_{\mathrm{res}} = \sum_S \Res_{\partial_S}(G_S \wedge \mu_S)$.

For $(\star)$: the de~Rham differential commutes with the residue
operation because $G$ is closed ($dG = 0$), so the cross-term
$[d_{\mathrm{dR}}, d_{\mathrm{res}}]$ reduces to
$\sum_S \Res_{\partial_S}(dG_S \wedge \mu_S) = 0$.

For $(\star\star)$: the square $d_{\mathrm{res}}^2$ involves
iterated residues along pairs of boundary strata.  For three
distinct points $i, j, k$, the iterated residue
$\Res_{\partial_{\{i,j\}}} \circ \Res_{\partial_{\{j,k\}}}$
plus cyclic permutations gives a sum of three terms.  The
Totaro relation
\[
\alpha_{ij} \cdot \alpha_{jk}
+ \alpha_{jk} \cdot \alpha_{ki}
+ \alpha_{ki} \cdot \alpha_{ij} \;=\; 0
\qquad
\text{in } H^{2(n-1)}(\Conf_3(\bR^n))
\]
at $n = 2$ ensures that this sum vanishes.  More precisely: each
term in $d_{\mathrm{res}}^2$ corresponds to a codimension-$2$
boundary stratum of $\overline{\Conf}_k(\bR^2)$, and the Totaro
relation is exactly the statement that the codimension-$2$
contributions cancel in cohomology.

For nested collisions (e.g., $\{i, j\} \subset \{i, j, k\}$),
the iterated residue factors through the operadic composition
$\mu_{\{i,j,k\}} = \mu_{\{i,j\}} \circ \mu_{\{j,k\}}$, and
cancellation follows from the $\Etwo$-algebra axiom (associativity
up to coherent homotopy, encoded by the FM operad structure).

For the algebraic contribution: iterated residues at the same
stratum $S$ produce the square $\mu_S \circ \mu_S$, which vanishes
by the $A_\infty$ relations encoded in the $\Etwo$ structure.

The argument generalizes Idrissi's result~\cite{Idrissi22} that
Poincar\'e duality on configuration spaces ensures $d^2 = 0$ for
the propagator-based differential, specialized to $n = 2$ with
explicit Totaro generators.
\end{proof}

\begin{remark}[Comparison with the chiral bar complex]
\label{rem:e2-vs-chiral}
\index{bar complex!$\Etwo$ vs chiral}
The $\Etwo$ bar complex of Definition~\ref{def:e2-bar-complex}
and the chiral bar complex of
Chapter~\ref{chap:bar-cobar} are related but distinct:
\begin{enumerate}
\item The \emph{chiral} bar complex uses the \emph{holomorphic}
  propagator $\omega_{ij} = d\log(z_i - z_j)$ on a complex curve~$X$,
  with residues extracting OPE poles.  It depends on the complex
  structure of~$X$.
\item The \emph{$\Etwo$} bar complex uses the \emph{topological}
  propagator $G = (1/2\pi)\, d\arg(z_1 - z_2)$ on a real surface,
  with linking-sphere integrals extracting $\Etwo$-multiplication data.
  It depends only on the oriented topology of~$M$.
\item For a chiral algebra $\cA$ on a complex curve~$X$, the
  underlying $\Etwo$-algebra $A = \cA^{\mathrm{top}}$ (obtained by
  forgetting holomorphic structure) satisfies
  $\barB_{\Etwo}(A) \simeq \barB^{\mathrm{ch}}(\cA)^{\mathrm{top}}$:
  the $\Etwo$ bar complex is the topological shadow of the chiral one.
\item The chiral bar complex carries strictly more information: the
  holomorphic propagator distinguishes conformal structures, while the
  $\Etwo$ propagator does not.  The genus-$g$ quantum corrections
  $\kappa(\cA) \cdot \lambda_g$ (Theorem~\ref{thm:genus-universality})
  live in the chiral bar complex but not in the $\Etwo$ bar complex
  (which sees only the underlying topological surface).
\end{enumerate}
\end{remark}


%================================================================
% SECTION 4: E_n KOSZUL DUALITY THEOREM
%================================================================

\section{$\En$ Koszul duality}
\label{sec:en-koszul-duality}
\index{Koszul duality!$\En$|textbf}

\begin{definition}[$\En$ bar-cobar adjunction]
\label{def:en-bar-cobar}
For each $n \geq 1$, the \emph{$\En$ bar functor} and \emph{$\En$
cobar functor} are defined by:
\begin{align}
\barB_{\En} &\colon
  \En\text{-}\mathrm{Alg} \;\longrightarrow\;
  \En\text{-}\mathrm{CoAlg},
  \label{eq:en-bar}\\
\Omega_{\En} &\colon
  \En\text{-}\mathrm{CoAlg} \;\longrightarrow\;
  \En\text{-}\mathrm{Alg},
  \label{eq:en-cobar}
\end{align}
where $\En\text{-}\mathrm{Alg}$ denotes the $\infty$-category of
$\En$-algebras in chain complexes and $\En\text{-}\mathrm{CoAlg}$
denotes the $\infty$-category of $\En$-coalgebras.  The bar
construction uses the chain complex
$C_*(\overline{\Conf}_k(\bR^n))$ with the differential of
\eqref{eq:e2-bar-diff} generalized to dimension~$n$, and the cobar
construction is its linear dual.
\end{definition}

\begin{theorem}[$\En$ Koszul duality;
\ClaimStatusProvedElsewhere]
\label{thm:en-koszul-duality}
\index{Koszul duality!$\En$}
\begin{enumerate}[label=\textup{(\roman*)}]
\item \textup{(Lurie~\cite{HA}, \S5.2)}\quad
  The functors $\barB_{\En}$ and $\Omega_{\En}$ form an adjoint pair
  \[
  \barB_{\En}
  \;\dashv\;
  \Omega_{\En}
  \colon
  \En\text{-}\mathrm{CoAlg}
  \;\rightleftarrows\;
  \En\text{-}\mathrm{Alg}.
  \]
\item \textup{(Francis~\cite{Francis2013}, Theorem~4.20)}\quad
  For an $\En$-algebra $A$ that is \emph{$\En$-Koszul}
  \textup{(}the $\En$ bar spectral sequence collapses\textup{)},
  the counit
  $\Omega_{\En}(\barB_{\En}(A)) \xrightarrow{\sim} A$
  is a quasi-isomorphism.
\item \textup{(Ayala--Francis~\cite{AF15}, Theorem~7.8)}\quad
  For a closed oriented $n$-manifold~$M$ and an $\En$-algebra~$A$,
  Poincar\'e--Koszul duality gives a natural equivalence:
  \begin{equation}\label{eq:poincare-koszul}
  \int_M A
  \;\simeq\;
  \int_M A^\vee[-n],
  \end{equation}
  where $\int_M$ denotes $\En$-factorization homology and $A^\vee$
  is the $\En$-Koszul dual coalgebra.
\end{enumerate}
\end{theorem}

\begin{remark}[Chain-level content of this monograph]
\label{rem:chain-level-en}
Part~(iii) of Theorem~\ref{thm:en-koszul-duality} is proved by
Ayala--Francis at the $\infty$-categorical level.  Our contribution
(at $n = 2$) is the \emph{chain-level} construction: explicit
propagators, explicit residue formulas, explicit spectral sequences.
Main Theorem~A (Theorem~\ref{thm:bar-cobar-isomorphism-main}) is
the chain-level refinement of part~(iii) at $n = 2$.  Main Theorem~B
(Theorem~\ref{thm:higher-genus-inversion}) is the analogue of
part~(ii) in the chiral setting, including the genus-$g$ quantum
corrections that have no counterpart in the purely topological $\En$
framework.  Main Theorem~C
(Theorem~\ref{thm:quantum-complementarity-main}) is the
complementarity of quantum corrections---a phenomenon specific to
curves ($n = 2$) with complex structure, invisible to the topological
$\Etwo$ theory.
\end{remark}

\begin{corollary}[Recovery of chiral bar-cobar at $n = 2$;
\ClaimStatusProvedHere]
\label{cor:n2-recovery}
For a chiral algebra $\cA$ on a complex curve~$X$, the chiral
bar-cobar adjunction of
Theorem~\textup{\ref{thm:bar-cobar-isomorphism-main}} is recovered
from the $\Etwo$ bar-cobar adjunction
\textup{(}Theorem~\textup{\ref{thm:en-koszul-duality}(i))} by
equipping the underlying $\Etwo$-algebra $\cA^{\mathrm{top}}$ with
the holomorphic propagator $\omega_{ij} = d\log(z_i - z_j)$ in
place of the topological propagator $G$.
\end{corollary}

\begin{proof}
The holomorphic propagator $\omega_{ij}$ on a curve~$X$ is a specific
representative of the $\Etwo$ propagator class
$[\alpha_{ij}] \in H^1(\Conf_2(X))$; different representatives give
quasi-isomorphic bar complexes (since they differ by exact forms,
and the bar differential is insensitive to exact corrections in
cohomology).  The Arnold relations on~$X$ are the Totaro relations at
$n = 2$, so $d^2 = 0$ follows from
Theorem~\ref{thm:e2-d-squared}.  The residue operation
$\Res_{z_i = z_j}$ is the $n = 2$ linking-sphere residue
(Proposition~\ref{prop:linking-sphere-residue}): the linking sphere
$S^1$ around the diagonal in a surface is the contour of integration
in the OPE residue formula.
\end{proof}


%================================================================
% SECTION 5: CONNECTION TO AYALA--FRANCIS
%================================================================

\section{Connection to Ayala--Francis}
\label{sec:ayala-francis-connection}
\index{Ayala--Francis!chain-level refinement}
\index{Poincar\'e--Koszul duality!chain-level}

Ayala--Francis~\cite{AF15} establish Poincar\'e--Koszul duality as
an equivalence of $\infty$-categories:

\begin{theorem}[Poincar\'e--Koszul duality, AF {\cite{AF15}};
\ClaimStatusProvedElsewhere]
\label{thm:af-pkd}
\index{Poincar\'e--Koszul duality}
Let $M$ be a closed oriented $n$-manifold and $A$ an $\En$-algebra.
There is a natural equivalence
\[
\int_M A
\;\simeq\;
\mathrm{Map}_{\En\text{-}\mathrm{CoAlg}}
\bigl(\barB_{\En}(\mathbf{1}),\, \barB_{\En}(A)\bigr),
\]
where $\mathbf{1}$ is the unit $\En$-algebra.
For $A$ augmented, this simplifies to the Poincar\'e--Koszul duality
\eqref{eq:poincare-koszul}.
\end{theorem}

\begin{proposition}[Our construction refines AF at $n = 2$;
\ClaimStatusProvedHere]
\label{prop:refines-af}
At $n = 2$, the chain-level bar complex
$\barB^{\mathrm{ch}}(\cA)$ of
Chapter~\textup{\ref{chap:bar-cobar}} is a
strict model for the $\Etwo$ bar construction of
Theorem~\textup{\ref{thm:af-pkd}}, in the following sense:
\begin{enumerate}[label=\textup{(\alph*)}]
\item The FM propagator integrals
  $\int_{\overline{C}_k(X)} \prod_{(i,j)} \omega_{ij}
  \otimes \prod_v a_v$
  compute the factorization homology $\int_X \cA$ as a chain
  complex.
\item The Verdier duality pairing on $\overline{C}_k(X)$
  \textup{(}\S\textup{\ref{sec:verdier-ayala-francis})} realizes the
  Poincar\'e duality of Theorem~\textup{\ref{thm:af-pkd}} at the
  cochain level.
\item The spectral sequence from bar filtration to bar cohomology
  is the descent spectral sequence of factorization homology with
  respect to the FM cover.
\end{enumerate}
\end{proposition}

\begin{proof}
Part (a): by the fundamental theorem of factorization
homology~\cite{AF15}, $\int_X \cA$ is computed by a \v{C}ech complex
with respect to any Weiss cover of~$X$ by disks.  The FM
compactification $\overline{C}_k(X)$ provides a canonical
resolution: configuration space integrals against the propagator
$\omega_{ij}$ compute the \v{C}ech differential, and the bar
elements $\bigotimes_v a_v$ provide the coefficients.  This is the
content of the Prism Principle
(Theorem~\ref{thm:prism-higher-genus}).

Part (b): Verdier duality $\mathbb{D}_{\overline{C}_k(X)}$ exchanges the
FM compactification with the open configuration space.  This is the
geometric incarnation of the Poincar\'e duality in
Theorem~\ref{thm:af-pkd}: the duality pairing between
$\int_X A$ and $\int_X A^\vee$ is realized by the intersection
pairing on $\overline{C}_k(X)$.  See
Theorem~\ref{thm:verdier-bar-cobar} for the precise statement
at $n = 2$.

Part (c): the FM cover filtration
$F_p = \bigoplus_{k \leq p} C_*(\overline{C}_k(X))
\otimes (s^{-1} \bar{A})^{\otimes k}$
gives a spectral sequence with
$E_1^{p,q} = H^q(\overline{C}_p(X)) \otimes (s^{-1}\bar{A})^{\otimes p}$
converging to $H^*(\int_X A)$.  This is the bar spectral sequence
of Theorem~\ref{thm:bar-cobar-spectral-sequence}, reinterpreted as
the descent spectral sequence for the factorization homology
presheaf.
\end{proof}


%================================================================
% SECTION 6: THE GENERAL n CASE AND FORMALITY
%================================================================

\section{General $n$ and Kontsevich formality}
\label{sec:general-n-formality}
\index{formality!$\En$ operad}

\subsection{The $\En$ bar complex for arbitrary $n$}

\begin{definition}[$\En$ bar complex, general $n$]
\label{def:en-bar-general}
\index{bar complex!general $n$}
For $n \geq 1$ and an $\En$-algebra $A$, the \emph{$\En$ bar complex}
is:
\begin{equation}\label{eq:en-bar-general}
\barB_{\En}(A)
\;=\;
\bigoplus_{k \geq 1}
\Omega^*\bigl(\overline{\Conf}_k(\bR^n)\bigr)
\otimes_{\Sigma_k} (s^{-1}\bar{A})^{\otimes k},
\quad
d \;=\; d_{\mathrm{dR}} + d_{\mathrm{res}},
\end{equation}
where $d_{\mathrm{res}} = \sum_S \Res_{\partial_S}(G_S \wedge \mu_S)$
uses the $\En$ propagator of Definition~\textup{\ref{def:en-propagator}}
and linking-sphere residues
\textup{(}Proposition~\textup{\ref{prop:linking-sphere-residue})}.
\end{definition}

\begin{theorem}[$d^2 = 0$ for the $\En$ bar complex;
\ClaimStatusProvedElsewhere]
\label{thm:en-d-squared}
\index{bar complex!$d^2 = 0$}
For any $n \geq 1$ and any $\En$-algebra $A$, the $\En$ bar
differential satisfies $d^2 = 0$.
\end{theorem}

\begin{proof}
The argument is identical to the $n = 2$ case
(Theorem~\ref{thm:e2-d-squared}), with the Arnold/Totaro
relations in $H^*(\Conf_3(\bR^n))$ providing the cancellation of
$d_{\mathrm{res}}^2$.  For general~$n$, the Totaro relation holds
in degree $2(n-1)$, and the cancellation works in any degree.

The result also follows from the general formality framework of
Idrissi~\cite{Idrissi22}: Poincar\'e duality on the FM
compactification $\overline{\Conf}_k(\bR^n)$ provides an explicit
quasi-isomorphism between the de~Rham complex of
$\overline{\Conf}_k(\bR^n)$ and its cohomology algebra (the Totaro
algebra), and the propagator differential is compatible with this
quasi-isomorphism.
\end{proof}

\begin{remark}[The dimensional ladder]
\label{rem:dimensional-ladder}
\index{dimensional ladder}
The transition from $n$ to $n+1$ trades algebraic structure for
topological generality:
\begin{center}
\begin{tabular}{clll}
$n$ & Propagator & Relations & Bar complex computes \\
\hline
$1$ & $d\log|t_1 - t_2|$ & Stasheff polytopes & Hochschild cohomology \\
$2$ & $d\log(z_1 - z_2)$ & Arnold (OS) & Chiral Koszul dual \\
$3$ & linking $2$-form & graph cohomology & Vassiliev invariants \\
$n$ & $(n{-}1)$-form & Totaro & $\En$-factorization homology
\end{tabular}
\end{center}
At each level, the propagator is a closed $(n-1)$-form, the
relations in $H^*(\Conf_3(\bR^n))$ ensure $d^2 = 0$, and the bar
cohomology computes the Koszul dual in the appropriate operadic
sense.
\end{remark}


\subsection{Knudsen's higher enveloping algebras}

\begin{theorem}[Higher enveloping algebras; \ClaimStatusProvedElsewhere]
\label{thm:knudsen-higher-enveloping}
\index{enveloping algebra!higher}
\textup{(}Knudsen~\cite{Knudsen18}\textup{)}\quad
For a Lie algebra $\fg$ and $n \geq 1$, the \emph{$n$-th enveloping
algebra} $U_n(\fg)$ is an $\En$-algebra satisfying:
\begin{enumerate}[label=\textup{(\roman*)}]
\item $U_1(\fg) = U(\fg)$ is the classical universal enveloping
  algebra.
\item \textup{(PBW)}\quad There is an equivalence of graded objects
  $\mathrm{gr}\, U_n(\fg) \simeq \mathrm{Sym}^{\En}(\fg[n-1])$,
  where $\mathrm{Sym}^{\En}$ is the free commutative $\En$-algebra.
\item \textup{(Koszul duality)}\quad
  $\barB_{\En}(U_n(\fg))$ is quasi-isomorphic to
  $\mathrm{CE}^{\En}_*(\fg)$, the $\En$-Chevalley--Eilenberg
  complex of~$\fg$.
\end{enumerate}
\end{theorem}

This result is the higher-dimensional analogue of the PBW theorem
and the classical bar-CE comparison.  At $n = 2$, the $\Etwo$
enveloping algebra $U_2(\fg)$ is the topological shadow of the
chiral universal enveloping algebra
$\widehat{U}(\fg_k) = V_k(\fg)$ (the vacuum module).  Knudsen's
Koszul duality at $n = 2$ is the topological shadow of our chiral
Koszul duality.


\subsection{Kontsevich formality and deformation quantization}

\begin{theorem}[Formality of $\Etwo$; \ClaimStatusProvedElsewhere]
\label{thm:e2-formality}
\index{formality!$\Etwo$ operad}
\textup{(}Kontsevich~\cite{Kontsevich03},
Tamarkin~\cite{Tamarkin00}\textup{)}\quad
The $\Etwo$ operad is formal over $\bQ$: there is a
quasi-isomorphism of operads
\[
C_*(\Etwo;\, \bQ)
\;\simeq\;
H_*(\Etwo;\, \bQ).
\]
Consequently, every $\Etwo$-algebra structure on a
cochain complex $A$ is determined up to quasi-isomorphism by the
induced $H_*(\Etwo)$-algebra structure on cohomology.
\end{theorem}

\begin{remark}[Formality and deformation quantization]
\label{rem:formality-dq}
\index{deformation quantization!from $\Etwo$ formality}
Kontsevich's formality theorem~\cite{Kontsevich03} proves that
$C^*(\bR^n, \cO)$ is formal as an $\Etwo$-algebra, which implies
the existence of star products on arbitrary Poisson manifolds.
In our framework
(Chapter~\ref{chap:chiral-deformation}), this is the
statement that the $\Etwo$ bar complex of the algebra of
polyvector fields is quasi-isomorphic to the $\Etwo$ bar complex
of the algebra of polydifferential operators---i.e., Koszul
duality at the $\Etwo$ level exchanges Poisson brackets with
star products.  The explicit formulas for the Kontsevich star
product arise from the propagator integrals on
$\overline{\Conf}_k(\mathcal{H})$ (the upper half-plane
compactification), which is a specific model for the $\Etwo$
bar complex with boundary conditions.
\end{remark}

\begin{proposition}[$\En$ formality for $n \geq 2$;
\ClaimStatusProvedElsewhere]
\label{prop:en-formality}
\index{formality!$\En$ operad}
\textup{(}Fresse--Willwacher, Idrissi~\cite{Idrissi22}\textup{)}\quad
The $\En$ operad is formal over $\bR$ for all $n \geq 2$.
The formality map is constructed explicitly via configuration
space integrals on $\overline{\Conf}_k(\bR^n)$ using the propagator
of Definition~\textup{\ref{def:en-propagator}}, with the Kontsevich
graph expansion providing the explicit formulas.
\end{proposition}


%================================================================
% SECTION 7: THE n = 3 CASE AND CHERN--SIMONS
%================================================================

\section{The $n = 3$ case: Chern--Simons theory}
\label{sec:n3-chern-simons}
\index{Chern--Simons!$\mathsf{E}_3$ bar complex}

The $n = 3$ case of $\En$ Koszul duality connects directly to
perturbative Chern--Simons theory and the Kontsevich integral
of Chapter~\ref{ch:kontsevich-integral}.

\begin{conjecture}[$\mathsf{E}_3$ bar complex and Chern--Simons;
\ClaimStatusConjectured]
\label{conj:e3-chern-simons}
\index{Chern--Simons!perturbative expansion}
For $\fg$ a finite-dimensional Lie algebra with invariant inner
product, the $\mathsf{E}_3$ bar complex
$\barB_{\mathsf{E}_3}(U_3(\fg))$ computes the perturbative
Chern--Simons invariants of $3$-manifolds:
\begin{enumerate}[label=\textup{(\roman*)}]
\item The propagator is the Chern--Simons propagator of
  Example~\textup{\ref{ex:n3-propagator}}: the Green's $2$-form
  of the gauge-fixed Laplacian on~$M$.
\item The Feynman diagrams of Axelrod--Singer are the
  contributions from individual boundary strata of
  $\overline{\Conf}_k(M)$.
\item The bar cohomology
  $H^*(\barB_{\mathsf{E}_3}(U_3(\fg)))$
  should recover the Kontsevich graph complex with Lie algebra
  coefficients.
\end{enumerate}
\end{conjecture}

\begin{remark}[Scope]
\label{rem:e3-scope}
The identification of Feynman diagrams with boundary strata
(part (ii)) is well-established in the literature
(Axelrod--Singer, Kontsevich~\cite{Kon94}, Cattaneo--Mnev~\cite{CattaneoMnev10}).  What is conjectured is
the full chain-level identification (iii), which would require
showing that the Axelrod--Singer compactification provides a
model for the $\mathsf{E}_3$ bar complex, not just individual
Feynman integrals.  The $n = 2$ analogue of this identification
is our Prism Principle
(Theorem~\ref{thm:prism-higher-genus}), which identifies the
chiral bar complex with the Feynman transform of the modular
operad.  The $n = 3$ case requires the theory of \emph{modular
$\mathsf{E}_3$-operads}, which is not yet developed.
\end{remark}


%================================================================
% SECTION 8: SUMMARY
%================================================================

\section{Summary and the dimensional ladder}
\label{sec:en-summary}

\noindent
\textbf{What is proved.}
\begin{enumerate}
\item The Totaro presentation of $H^*(\Conf_k(\bR^n))$ generalizes
  Arnold's presentation from $n = 2$ to all~$n$
  (Theorem~\ref{thm:totaro-presentation}).
\item The $\Etwo$ bar complex satisfies $d^2 = 0$ via the Totaro
  relations (Theorem~\ref{thm:e2-d-squared}).
\item The chiral bar-cobar duality of Main Theorem~A is recovered
  as the $n = 2$ case of $\En$ Koszul duality
  (Corollary~\ref{cor:n2-recovery}).
\item Our construction is the chain-level refinement of
  Ayala--Francis Poincar\'e--Koszul duality at $n = 2$
  (Proposition~\ref{prop:refines-af}).
\item $\En$ Koszul duality, $\En$ formality, and higher enveloping
  algebras are established in the literature for all~$n$
  (Theorem~\ref{thm:en-koszul-duality},
  Theorem~\ref{thm:knudsen-higher-enveloping},
  Proposition~\ref{prop:en-formality}).
\end{enumerate}

\noindent
\textbf{What remains.}
\begin{enumerate}
\item Explicit chain-level formulas for the $\En$ bar complex at
  $n \geq 3$ (the propagator and FM compactification exist, but
  the full chain-level identification with Feynman diagrams
  requires modular $\En$-operads).
\item The $n = 3$ Chern--Simons identification
  (Conjecture~\ref{conj:e3-chern-simons}): connecting the
  $\mathsf{E}_3$ bar complex to perturbative CS invariants at the
  chain level.
\item Higher-genus $\En$ theory: for $n \geq 3$, the analogue of
  genus-$g$ quantum corrections would require a notion of
  ``higher-dimensional moduli'' replacing $\overline{\mathcal{M}}_g$.
  The correct framework may be cobordism categories or extended
  TQFTs, but this is far from the current state of the art.
\end{enumerate}

The dimensional ladder makes precise the sense in which this
monograph's results---chiral bar-cobar duality, quantum
complementarity, genus expansions---are the $n = 2$ chapter
of a general story that extends through all dimensions.
